% --- Archon physics formalization source begin ---
% archon:physics
% archon:covers PhyXMiniProblems/problem_phyx_mini_0580.lean
% archon:source-report reports/phyx_mini/problem_phyx_mini_0580.source.json

\chapter{Physics problem phyx\_mini\_0580}
\label{ch:PhyXMiniProblems_problem_phyx_mini_0580}

\paragraph{Problem source.}
\#\# Physical scenario

Figure is an energy-level diagram for a fictitious three-dimensional infinite potential well that contains one electron. The number of degenerate states of the levels are indicated: "non" means nondegenerate (which includes the ground state) and "triple" means 3 states. Assume that the electrostatic forces between the electrons are negligible.

\#\# Question

If we put a total of 22 electrons in the well, what multiple of \textbackslash{}( h\textasciicircum{}2 \textbackslash{}! / \textbackslash{}! 8mL\textasciicircum{}2 \textbackslash{}) gives the energy of the ground state of the 22-electron system?

\#\# Answer choices

- A: A: 174\textbackslash{}left(\textbackslash{}frac\{h\textasciicircum{}2\}\{8mL\textasciicircum{}2\}\textbackslash{}right)

- B: B: 192\textbackslash{}left(\textbackslash{}frac\{h\textasciicircum{}2\}\{8mL\textasciicircum{}2\}\textbackslash{}right)

- C: C: 186\textbackslash{}left(\textbackslash{}frac\{h\textasciicircum{}2\}\{8mL\textasciicircum{}2\}\textbackslash{}right

- D: D: 162\textbackslash{}left(\textbackslash{}frac\{h\textasciicircum{}2\}\{8mL\textasciicircum{}2\}\textbackslash{}right)

\#\# Figure caption (auxiliary; use the image as primary evidence)

The image is a diagram related to energy levels. It displays a vertical axis labeled \textbackslash{}( E \textbackslash{}left(\textbackslash{}frac\{h\textasciicircum{}2\}\{8mL\textasciicircum{}2\}\textbackslash{}right) \textbackslash{}) on the left and horizontal lines corresponding to various energy levels. The levels are numbered as 14, 12, 11, 9, 6, and 3 from top to bottom. Each level is associated with a descriptive label on the right:

- 14 is labeled "Triple"
- 12 is labeled "Non"
- 11 is labeled "Triple"
- 9 is labeled "Triple"
- 6 is labeled "Triple"
- 3 is labeled "Ground"

The lines are horizontal and colored blue, indicating different energy states.

\#\# Dataset metadata

Quantum Phenomena; Physical Model Grounding Reasoning, Multi-Formula Reasoning

\paragraph{Recorded answer/context.}
C: C: 186\textbackslash{}left(\textbackslash{}frac\{h\textasciicircum{}2\}\{8mL\textasciicircum{}2\}\textbackslash{}right

\paragraph{Figure/image path.}
/root/proposal\_for\_physic/science-mango/phyx\_mini\_run/phyx\_data/test\_image/580.png

\paragraph{Formalization target.}
create a compiling Lean file with sorry bodies at `PhyXMiniProblems/problem\_phyx\_mini\_0580.lean`. The Lean declarations must preserve the physical quantities, dimensions or dimensional roles, figure labels, governing-law hypotheses, and final relation expressed by this problem.
Use Mathlib/Physlib names found through LeanExplore where available. If a physics API is missing, introduce faithful local abstractions rather than scalar placeholder aliases.

\begin{theorem}[Physics formalization target]
\label{thm:physics:phyx_mini_0580:target}
\lean{PhyXMiniProblems.ProblemPhyXMini0580.problem_phyx_mini_0580}
\uses{def:physics:phyx-mini-0580:phyxminiproblems-problemphyxmini0580-energyinjoules, def:physics:phyx-mini-0580:phyxminiproblems-problemphyxmini0580-manyelectroninfinitewellsetup, def:physics:phyx-mini-0580:phyxminiproblems-problemphyxmini0580-infinitewellenergyscaleinjoules, def:physics:phyx-mini-0580:phyxminiproblems-problemphyxmini0580-matchesinfinitewellscenario, def:physics:phyx-mini-0580:phyxminiproblems-problemphyxmini0580-matchesinfinitewellfigure, def:physics:phyx-mini-0580:phyxminiproblems-problemphyxmini0580-hasphysicalinfinitewellparameters, def:physics:phyx-mini-0580:phyxminiproblems-problemphyxmini0580-satisfiesindependentelectronfillinglaws, def:physics:phyx-mini-0580:phyxminiproblems-problemphyxmini0580-displayedgroundenergymultiplier, def:physics:phyx-mini-0580:phyxminiproblems-problemphyxmini0580-recordedanswerchoice}
The assigned autoformalize agent should translate this physics problem into Lean declarations in the covered file, with theorem and lemma proof bodies written as `by sorry`.
\end{theorem}
\begin{proof}
This is an autoformalization task, not a proof task. Produce statements that can later be proved by the physics prover without weakening the physical model.
\end{proof}

% --- Archon declaration topology begin ---
\section{Declaration topology}
\begin{definition}[Length Quantity]
  \label{def:physics:phyx-mini-0580:phyxminiproblems-problemphyxmini0580-lengthquantity}
  \lean{PhyXMiniProblems.ProblemPhyXMini0580.LengthQuantity}
  A nonnegative, unit-independent physical length
\end{definition}
\begin{proof}
  This is the declared model definition of Length Quantity.
\end{proof}
\begin{definition}[Mass Quantity]
  \label{def:physics:phyx-mini-0580:phyxminiproblems-problemphyxmini0580-massquantity}
  \lean{PhyXMiniProblems.ProblemPhyXMini0580.MassQuantity}
  A nonnegative, unit-independent physical mass
\end{definition}
\begin{proof}
  This is the declared model definition of Mass Quantity.
\end{proof}
\begin{definition}[action Dimension]
  \label{def:physics:phyx-mini-0580:phyxminiproblems-problemphyxmini0580-actiondimension}
  \lean{PhyXMiniProblems.ProblemPhyXMini0580.actionDimension}
  The physical dimension M L² T⁻¹ of action
\end{definition}
\begin{proof}
  This is the declared model definition of action Dimension.
\end{proof}
\begin{definition}[Action Quantity]
  \label{def:physics:phyx-mini-0580:phyxminiproblems-problemphyxmini0580-actionquantity}
  \lean{PhyXMiniProblems.ProblemPhyXMini0580.ActionQuantity}
  \uses{def:physics:phyx-mini-0580:phyxminiproblems-problemphyxmini0580-actiondimension}
  A nonnegative physical action, used here for the ordinary Planck constant
\end{definition}
\begin{proof}
  This is the declared model definition of Action Quantity.
\end{proof}
\begin{definition}[length In Meters]
  \label{def:physics:phyx-mini-0580:phyxminiproblems-problemphyxmini0580-lengthinmeters}
  \lean{PhyXMiniProblems.ProblemPhyXMini0580.lengthInMeters}
  \uses{def:physics:phyx-mini-0580:phyxminiproblems-problemphyxmini0580-lengthquantity}
  Coherent-SI readout of a physical length, in metres
\end{definition}
\begin{proof}
  This is the declared model definition of length In Meters.
\end{proof}
\begin{definition}[mass In Kilograms]
  \label{def:physics:phyx-mini-0580:phyxminiproblems-problemphyxmini0580-massinkilograms}
  \lean{PhyXMiniProblems.ProblemPhyXMini0580.massInKilograms}
  \uses{def:physics:phyx-mini-0580:phyxminiproblems-problemphyxmini0580-massquantity}
  Coherent-SI readout of a physical mass, in kilograms
\end{definition}
\begin{proof}
  This is the declared model definition of mass In Kilograms.
\end{proof}
\begin{definition}[action In Joule Seconds]
  \label{def:physics:phyx-mini-0580:phyxminiproblems-problemphyxmini0580-actioninjouleseconds}
  \lean{PhyXMiniProblems.ProblemPhyXMini0580.actionInJouleSeconds}
  \uses{def:physics:phyx-mini-0580:phyxminiproblems-problemphyxmini0580-actionquantity}
  Coherent-SI readout of an action, in joule-seconds
\end{definition}
\begin{proof}
  This is the declared model definition of action In Joule Seconds.
\end{proof}
\begin{definition}[energy In Joules]
  \label{def:physics:phyx-mini-0580:phyxminiproblems-problemphyxmini0580-energyinjoules}
  \lean{PhyXMiniProblems.ProblemPhyXMini0580.energyInJoules}
  Coherent-SI readout of a Physlib energy, in joules
\end{definition}
\begin{proof}
  This is the declared model definition of energy In Joules.
\end{proof}
\begin{definition}[Infinite Well Level]
  \label{def:physics:phyx-mini-0580:phyxminiproblems-problemphyxmini0580-infinitewelllevel}
  \lean{PhyXMiniProblems.ProblemPhyXMini0580.InfiniteWellLevel}
  The six energy levels visible in the supplied diagram
\end{definition}
\begin{proof}
  This is the declared model definition of Infinite Well Level.
\end{proof}
\begin{definition}[Electron Spin]
  \label{def:physics:phyx-mini-0580:phyxminiproblems-problemphyxmini0580-electronspin}
  \lean{PhyXMiniProblems.ProblemPhyXMini0580.ElectronSpin}
  The two electron spin states that can occupy one spatial state
\end{definition}
\begin{proof}
  This is the declared model definition of Electron Spin.
\end{proof}
\begin{definition}[Degeneracy Label]
  \label{def:physics:phyx-mini-0580:phyxminiproblems-problemphyxmini0580-degeneracylabel}
  \lean{PhyXMiniProblems.ProblemPhyXMini0580.DegeneracyLabel}
  Literal degeneracy words printed next to the energy lines
\end{definition}
\begin{proof}
  This is the declared model definition of Degeneracy Label.
\end{proof}
\begin{definition}[spatial State Count]
  \label{def:physics:phyx-mini-0580:phyxminiproblems-problemphyxmini0580-degeneracylabel-spatialstatecount}
  \lean{PhyXMiniProblems.ProblemPhyXMini0580.DegeneracyLabel.spatialStateCount}
  \uses{def:physics:phyx-mini-0580:phyxminiproblems-problemphyxmini0580-degeneracylabel}
  Number of spatial one-electron states represented by a degeneracy word
\end{definition}
\begin{proof}
  This is the declared model definition of spatial State Count.
\end{proof}
\begin{definition}[Level Line Orientation]
  \label{def:physics:phyx-mini-0580:phyxminiproblems-problemphyxmini0580-levellineorientation}
  \lean{PhyXMiniProblems.ProblemPhyXMini0580.LevelLineOrientation}
  Orientation of a level line in the raster
\end{definition}
\begin{proof}
  This is the declared model definition of Level Line Orientation.
\end{proof}
\begin{definition}[Level Line Color]
  \label{def:physics:phyx-mini-0580:phyxminiproblems-problemphyxmini0580-levellinecolor}
  \lean{PhyXMiniProblems.ProblemPhyXMini0580.LevelLineColor}
  Color category of a level line in the raster
\end{definition}
\begin{proof}
  This is the declared model definition of Level Line Color.
\end{proof}
\begin{definition}[Vertical Axis Quantity]
  \label{def:physics:phyx-mini-0580:phyxminiproblems-problemphyxmini0580-verticalaxisquantity}
  \lean{PhyXMiniProblems.ProblemPhyXMini0580.VerticalAxisQuantity}
  Quantity named on the vertical axis
\end{definition}
\begin{proof}
  This is the declared model definition of Vertical Axis Quantity.
\end{proof}
\begin{definition}[Printed Energy Unit]
  \label{def:physics:phyx-mini-0580:phyxminiproblems-problemphyxmini0580-printedenergyunit}
  \lean{PhyXMiniProblems.ProblemPhyXMini0580.PrintedEnergyUnit}
  Symbolic unit expression printed beside the vertical-axis energy label
\end{definition}
\begin{proof}
  This is the declared model definition of Printed Energy Unit.
\end{proof}
\begin{definition}[Infinite Well Energy Figure]
  \label{def:physics:phyx-mini-0580:phyxminiproblems-problemphyxmini0580-infinitewellenergyfigure}
  \lean{PhyXMiniProblems.ProblemPhyXMini0580.InfiniteWellEnergyFigure}
  \uses{def:physics:phyx-mini-0580:phyxminiproblems-problemphyxmini0580-infinitewelllevel, def:physics:phyx-mini-0580:phyxminiproblems-problemphyxmini0580-degeneracylabel, def:physics:phyx-mini-0580:phyxminiproblems-problemphyxmini0580-levellineorientation, def:physics:phyx-mini-0580:phyxminiproblems-problemphyxmini0580-levellinecolor, def:physics:phyx-mini-0580:phyxminiproblems-problemphyxmini0580-verticalaxisquantity, def:physics:phyx-mini-0580:phyxminiproblems-problemphyxmini0580-printedenergyunit}
  Literal and geometric information carried by the primary energy diagram
\end{definition}
\begin{proof}
  This is the declared model definition of Infinite Well Energy Figure.
\end{proof}
\begin{definition}[Confinement Model]
  \label{def:physics:phyx-mini-0580:phyxminiproblems-problemphyxmini0580-confinementmodel}
  \lean{PhyXMiniProblems.ProblemPhyXMini0580.ConfinementModel}
  The confinement model stated in the problem
\end{definition}
\begin{proof}
  This is the declared model definition of Confinement Model.
\end{proof}
\begin{definition}[Confined Particle Species]
  \label{def:physics:phyx-mini-0580:phyxminiproblems-problemphyxmini0580-confinedparticlespecies}
  \lean{PhyXMiniProblems.ProblemPhyXMini0580.ConfinedParticleSpecies}
  The particle species loaded into the well
\end{definition}
\begin{proof}
  This is the declared model definition of Confined Particle Species.
\end{proof}
\begin{definition}[Many Electron Infinite Well Setup]
  \label{def:physics:phyx-mini-0580:phyxminiproblems-problemphyxmini0580-manyelectroninfinitewellsetup}
  \lean{PhyXMiniProblems.ProblemPhyXMini0580.ManyElectronInfiniteWellSetup}
  \uses{def:physics:phyx-mini-0580:phyxminiproblems-problemphyxmini0580-lengthquantity, def:physics:phyx-mini-0580:phyxminiproblems-problemphyxmini0580-massquantity, def:physics:phyx-mini-0580:phyxminiproblems-problemphyxmini0580-actionquantity, def:physics:phyx-mini-0580:phyxminiproblems-problemphyxmini0580-infinitewelllevel, def:physics:phyx-mini-0580:phyxminiproblems-problemphyxmini0580-infinitewellenergyfigure, def:physics:phyx-mini-0580:phyxminiproblems-problemphyxmini0580-confinementmodel, def:physics:phyx-mini-0580:phyxminiproblems-problemphyxmini0580-confinedparticlespecies}
  The dimensionful well parameters, spectrum, occupation, and total energy. The fields are independent data; all physical relations between them are imposed below as scenario, readout, or governing-law premises
\end{definition}
\begin{proof}
  This is the declared model definition of Many Electron Infinite Well Setup.
\end{proof}
\begin{definition}[infinite Well Energy Scale In Joules]
  \label{def:physics:phyx-mini-0580:phyxminiproblems-problemphyxmini0580-infinitewellenergyscaleinjoules}
  \lean{PhyXMiniProblems.ProblemPhyXMini0580.infiniteWellEnergyScaleInJoules}
  \uses{def:physics:phyx-mini-0580:phyxminiproblems-problemphyxmini0580-lengthinmeters, def:physics:phyx-mini-0580:phyxminiproblems-problemphyxmini0580-massinkilograms, def:physics:phyx-mini-0580:phyxminiproblems-problemphyxmini0580-actioninjouleseconds, def:physics:phyx-mini-0580:phyxminiproblems-problemphyxmini0580-manyelectroninfinitewellsetup}
  SI value of the scale h²/(8mL²) named on the figure axis
\end{definition}
\begin{proof}
  This is the declared model definition of infinite Well Energy Scale In Joules.
\end{proof}
\begin{definition}[level Electron Capacity]
  \label{def:physics:phyx-mini-0580:phyxminiproblems-problemphyxmini0580-levelelectroncapacity}
  \lean{PhyXMiniProblems.ProblemPhyXMini0580.levelElectronCapacity}
  \uses{def:physics:phyx-mini-0580:phyxminiproblems-problemphyxmini0580-infinitewelllevel, def:physics:phyx-mini-0580:phyxminiproblems-problemphyxmini0580-electronspin, def:physics:phyx-mini-0580:phyxminiproblems-problemphyxmini0580-infinitewellenergyfigure}
  Maximum number of electrons at a pictured level under the Pauli rule
\end{definition}
\begin{proof}
  This is the declared model definition of level Electron Capacity.
\end{proof}
\begin{definition}[Is Pauli Allowed Occupation]
  \label{def:physics:phyx-mini-0580:phyxminiproblems-problemphyxmini0580-ispauliallowedoccupation}
  \lean{PhyXMiniProblems.ProblemPhyXMini0580.IsPauliAllowedOccupation}
  \uses{def:physics:phyx-mini-0580:phyxminiproblems-problemphyxmini0580-infinitewelllevel, def:physics:phyx-mini-0580:phyxminiproblems-problemphyxmini0580-manyelectroninfinitewellsetup, def:physics:phyx-mini-0580:phyxminiproblems-problemphyxmini0580-levelelectroncapacity}
  An occupation has 22-particle bookkeeping and respects every Pauli cap
\end{definition}
\begin{proof}
  This is the declared model definition of Is Pauli Allowed Occupation.
\end{proof}
\begin{definition}[occupation Energy In Joules]
  \label{def:physics:phyx-mini-0580:phyxminiproblems-problemphyxmini0580-occupationenergyinjoules}
  \lean{PhyXMiniProblems.ProblemPhyXMini0580.occupationEnergyInJoules}
  \uses{def:physics:phyx-mini-0580:phyxminiproblems-problemphyxmini0580-energyinjoules, def:physics:phyx-mini-0580:phyxminiproblems-problemphyxmini0580-infinitewelllevel, def:physics:phyx-mini-0580:phyxminiproblems-problemphyxmini0580-manyelectroninfinitewellsetup}
  Sum of occupied one-particle energy readouts for an occupation
\end{definition}
\begin{proof}
  This is the declared model definition of occupation Energy In Joules.
\end{proof}
\begin{definition}[Matches Infinite Well Scenario]
  \label{def:physics:phyx-mini-0580:phyxminiproblems-problemphyxmini0580-matchesinfinitewellscenario}
  \lean{PhyXMiniProblems.ProblemPhyXMini0580.MatchesInfiniteWellScenario}
  \uses{def:physics:phyx-mini-0580:phyxminiproblems-problemphyxmini0580-manyelectroninfinitewellsetup}
  Prose data specifying the well, particles, initial diagram, and new load
\end{definition}
\begin{proof}
  This is the declared model definition of Matches Infinite Well Scenario.
\end{proof}
\begin{definition}[Matches Infinite Well Figure]
  \label{def:physics:phyx-mini-0580:phyxminiproblems-problemphyxmini0580-matchesinfinitewellfigure}
  \lean{PhyXMiniProblems.ProblemPhyXMini0580.MatchesInfiniteWellFigure}
  \uses{def:physics:phyx-mini-0580:phyxminiproblems-problemphyxmini0580-manyelectroninfinitewellsetup}
  Exact data transcribed from the primary raster. This contains no occupation numbers, many-electron ground-state energy, or requested answer multiplier
\end{definition}
\begin{proof}
  This is the declared model definition of Matches Infinite Well Figure.
\end{proof}
\begin{definition}[Has Physical Infinite Well Parameters]
  \label{def:physics:phyx-mini-0580:phyxminiproblems-problemphyxmini0580-hasphysicalinfinitewellparameters}
  \lean{PhyXMiniProblems.ProblemPhyXMini0580.HasPhysicalInfiniteWellParameters}
  \uses{def:physics:phyx-mini-0580:phyxminiproblems-problemphyxmini0580-lengthinmeters, def:physics:phyx-mini-0580:phyxminiproblems-problemphyxmini0580-massinkilograms, def:physics:phyx-mini-0580:phyxminiproblems-problemphyxmini0580-actioninjouleseconds, def:physics:phyx-mini-0580:phyxminiproblems-problemphyxmini0580-energyinjoules, def:physics:phyx-mini-0580:phyxminiproblems-problemphyxmini0580-manyelectroninfinitewellsetup}
  Positivity conditions for the physical well and its energy scale
\end{definition}
\begin{proof}
  This is the declared model definition of Has Physical Infinite Well Parameters.
\end{proof}
\begin{definition}[Satisfies Independent Electron Filling Laws]
  \label{def:physics:phyx-mini-0580:phyxminiproblems-problemphyxmini0580-satisfiesindependentelectronfillinglaws}
  \lean{PhyXMiniProblems.ProblemPhyXMini0580.SatisfiesIndependentElectronFillingLaws}
  \uses{def:physics:phyx-mini-0580:phyxminiproblems-problemphyxmini0580-energyinjoules, def:physics:phyx-mini-0580:phyxminiproblems-problemphyxmini0580-manyelectroninfinitewellsetup, def:physics:phyx-mini-0580:phyxminiproblems-problemphyxmini0580-infinitewellenergyscaleinjoules, def:physics:phyx-mini-0580:phyxminiproblems-problemphyxmini0580-ispauliallowedoccupation, def:physics:phyx-mini-0580:phyxminiproblems-problemphyxmini0580-occupationenergyinjoules}
  Governing laws for the independent-electron filling calculation: the named energy scale is h²/(8mL²); each one-particle energy is its printed multiplier times that scale; the shown occupation is Pauli-allowed and minimizes the noninteracting sum; electron-electron electrostatic interactions contribute zero in the stated approximation; and the many-electron ground-state energy is the occupied one-particle sum plus that interaction contribution. No field states a level occupation or the numerical total 186
\end{definition}
\begin{proof}
  This is the declared model definition of Satisfies Independent Electron Filling Laws.
\end{proof}
\begin{lemma}[ground state occupation for twenty two electrons]
  \label{lem:physics:phyx-mini-0580:phyxminiproblems-problemphyxmini0580-ground-state-occupation-for-twenty-two-electrons}
  \lean{PhyXMiniProblems.ProblemPhyXMini0580.ground_state_occupation_for_twenty_two_electrons}
  \uses{def:physics:phyx-mini-0580:phyxminiproblems-problemphyxmini0580-manyelectroninfinitewellsetup, def:physics:phyx-mini-0580:phyxminiproblems-problemphyxmini0580-matchesinfinitewellscenario, def:physics:phyx-mini-0580:phyxminiproblems-problemphyxmini0580-matchesinfinitewellfigure, def:physics:phyx-mini-0580:phyxminiproblems-problemphyxmini0580-hasphysicalinfinitewellparameters, def:physics:phyx-mini-0580:phyxminiproblems-problemphyxmini0580-satisfiesindependentelectronfillinglaws}
  The Pauli caps and strict energy ordering force the first five levels to be filled as 2, 6, 6, 6, 2, leaving the level at multiplier 14 empty. This is a derived conclusion, not a premise of the main theorem
\end{lemma}
\begin{proof}
  The conclusion follows from the governing relations and figure readouts named in the statement of ground state occupation for twenty two electrons.
\end{proof}
\begin{definition}[Answer Choice]
  \label{def:physics:phyx-mini-0580:phyxminiproblems-problemphyxmini0580-answerchoice}
  \lean{PhyXMiniProblems.ProblemPhyXMini0580.AnswerChoice}
  Labels of the four displayed answers
\end{definition}
\begin{proof}
  This is the declared model definition of Answer Choice.
\end{proof}
\begin{definition}[displayed Ground Energy Multiplier]
  \label{def:physics:phyx-mini-0580:phyxminiproblems-problemphyxmini0580-displayedgroundenergymultiplier}
  \lean{PhyXMiniProblems.ProblemPhyXMini0580.displayedGroundEnergyMultiplier}
  \uses{def:physics:phyx-mini-0580:phyxminiproblems-problemphyxmini0580-answerchoice}
  Dimensionless energy multiplier printed beside each answer choice
\end{definition}
\begin{proof}
  This is the declared model definition of displayed Ground Energy Multiplier.
\end{proof}
\begin{definition}[recorded Answer Choice]
  \label{def:physics:phyx-mini-0580:phyxminiproblems-problemphyxmini0580-recordedanswerchoice}
  \lean{PhyXMiniProblems.ProblemPhyXMini0580.recordedAnswerChoice}
  \uses{def:physics:phyx-mini-0580:phyxminiproblems-problemphyxmini0580-answerchoice}
  Dataset metadata records answer choice C
\end{definition}
\begin{proof}
  This is the declared model definition of recorded Answer Choice.
\end{proof}
% --- Archon declaration topology end ---
% --- Archon physics formalization source end ---
